\section{Introduction}


Large language models excel at reasoning but require external tools to interact with the world. The Model Context Protocol (MCP) provides a standardized JSON-RPC interface for tool discovery (\texttt{tools/list}), invocation (\texttt{tools/call}), and resource access. However, existing MCP implementations offer limited tool coverage, forcing developers to implement custom tools for common operations.

Hanzo MCP Server provides a batteries-included implementation with tools spanning the full software development lifecycle. The server is designed for three deployment scenarios:
\begin{enumerate}
    \item \textbf{Local Development}: Direct integration with IDEs and coding assistants.
    \item \textbf{Remote Execution}: Sandboxed execution in cloud environments.
    \item \textbf{Agent Orchestration}: Tool backend for multi-agent systems.
\end{enumerate}

\paragraph{Contributions.} This paper makes the following contributions:
\begin{itemize}
    \item A configurable tool registry supporting category-based filtering and runtime tool addition/removal.
    \item A unified search engine combining ripgrep-based text search, tree-sitter AST analysis, and vector similarity retrieval.
    \item Native Playwright integration for cross-browser automation with device emulation.
    \item Comprehensive benchmarks across latency, throughput, and accuracy dimensions.
\end{itemize}

